\section{Perturbation tools and a completely positive criterion}
\label{sec:foundations}

Several steps of the proof produce an averaged map instead of a
homomorphism. The main result of this section, \Cref{prop:ucp}, turns
such a map into a small unitary conjugacy. All constants in this section
are independent of the algebras and Hilbert spaces.

\subsection{Metrics, amplification, and limits}

We use the unit-ball formula \eqref{eq:kk-definition} also for unital
$C^*$-algebras. For $A\subseteq\B(H)$, write $\U(A)$ for its unitary
group, $A'=\{T\in\B(H):Ta=aT\text{ for every }a\in A\}$ for its
commutant, and $Z(A)=A\cap A'$ for its center. The symbol
$A\barotimes B$ denotes the spatial von Neumann tensor product.

For a linear map $T:A\to\B(H)$, let $T_n$ act entrywise on $M_n(A)$,
and put $\norm{T}_{\cb}=\sup_n\norm{T_n}$.
The same notation for rectangular matrices is obtained by adding zero
rows and columns and compressing. For concretely represented von Neumann
algebras define
\[
 \dcb(A,B)=\sup_{n\ge1}d(M_n(A),M_n(B)).
\]
Here every matrix unit ball uses its operator norm on $H^n$.
We write $A\subset_\gamma B$ if there is $\gamma'<\gamma$ such that
every $x\in A_1$ has $\dist(x,B)\le\gamma'$; the approximant in $B$
need not be contractive. If $\norm{x}\le1$ and $\norm{x-y}\le r$, then
\[
 y_0=\frac{y}{\max\{1,\norm y\}}\in B_1,
 \qquad \norm{x-y_0}\le2r.
\]
We use this normalization whenever passing from near inclusions to
unit-ball distances. Conversely, a strict upper bound for $d(A,B)$
gives both near inclusions with that bound. For a unitary $v$,
\begin{equation}\label{eq:gap-adjustment}
 \begin{aligned}
 d(A,vBv^*)&\le d(A,B)+2\norm{v-I},\\
 \dcb(A,vBv^*)&\le\dcb(A,B)+2\norm{v-I}.
 \end{aligned}
\end{equation}
These inequalities follow at every matrix level from
$\norm{v_n x v_n^*-x}\le2\norm{v-I}\norm x$.
Also $\norm{v_1\cdots v_k-I}\le\sum_j\norm{v_j-I}$.

Inner products are linear in their first variable. A normal map between
von Neumann algebras is an ultraweakly continuous map. A map is
\emph{unital completely positive}, abbreviated UCP, if it preserves the
identity and each of its matrix amplifications is positive. Such maps
are completely contractive.

We repeatedly use the following elementary compactness facts. A bounded
operator ball is compact in the weak operator topology (WOT). A net of
UCP maps with fixed domain and codomain $\B(H)$ has a subnet converging pointwise in WOT:
use the product of the operator balls indexed by domain elements.
The limit is UCP, since linearity, the identity, and every finite matrix
positivity inequality pass to the limit. If the maps are within $a$ in
cb norm of a fixed map, so is their limit; apply WOT closedness of norm
balls at each fixed matrix size. If their values lie in a fixed von
Neumann algebra, the limit has the same property. Normality need not
pass to this limit and will not be required.

For clarity, $O(t)$ denotes a bound $Ct$ for an absolute constant $C$,
valid below an absolute positive threshold for $t$. Constants may change
between displays. A function tending to zero will always be chosen
uniformly over the explicitly stated class of algebras.

\subsection{The amenable embedding theorem}

The following form of Christensen's theorem supplies the small changes
of position used in the proof \cite{Christensen1980}.

\begin{theorem}[Amenable near inclusions]\label{thm:amenable}
Let $R,Q\subseteq\B(H)$ be von Neumann algebras with identity $I$, with
$R$ amenable. If $R\subset_\gamma Q$ and $\gamma<1/100$, there is a
unitary $v\in(R\cup Q)''$ such that
\[
 vRv^*\subseteq Q,\qquad \norm{v-I}\le150\gamma.
\]
If also $Q\subset_\gamma R$ and $\gamma<1/101$, the same conclusion
holds with $vRv^*=Q$.
\end{theorem}

This is \cite[Theorem~3.3(i)]{CCSSWW2014}, with its strict
near-inclusion convention. The location of $v$ in $(R\cup Q)''$ is
useful: any algebra commuting with both $R$ and $Q$ also commutes with
the change of position. We apply the theorem to abelian algebras,
hyperfinite finite factors, and type $\mathrm I$ algebras, all of which
are amenable.

\subsection{From commutators to intertwiners}

For $h\in\B(H)$ write $\ad_h(x)=[h,x]=hx-xh$.

\begin{lemma}\label{lem:commutant-cb}
For every unital von Neumann algebra $A\subseteq\B(H)$ and
$h\in\B(H)$,
\begin{equation}\label{eq:commutant-cb}
 \dist(h,A')\le\norm{\ad_h|_A}_{\cb}.
\end{equation}
Consequently $\dcb(A,B)\le\gamma$ implies $d(A',B')\le4\gamma$.
\end{lemma}

\begin{proof}
Since $A'$ is weak-star closed, its preannihilator in the trace class
computes the quotient norm $\dist(h,A')$. Let $\varphi$ be a trace-class
functional of norm at most one vanishing on $A'$. A singular-value
decomposition gives vectors $\xi,\eta\in\ell^2\otimes H$, each of
norm at most one, with
\[
 \varphi(x)=\ip{(I\otimes x)\xi}{\eta}\quad(x\in\B(H)).
\]
This uses only countably many vectors even when $H$ is nonseparable.
Let $p$ project onto $\overline{(I\otimes A')\xi}$. This subspace
reduces $I\otimes A'$, so every matrix entry of $p$ belongs to $A$.
Moreover $p\xi=\xi$ and $p\eta=0$. Therefore
\[
 \varphi(h)=\ip{[I\otimes h,p]\xi}{\eta}.
\]
If $q_n$ projects onto the first $n$ coordinates of $\ell^2$, the
compression of this commutator is
$[I_n\otimes h,q_npq_n]$, of norm at most
$\norm{\ad_h|_A}_{\cb}$ because $q_npq_n\in M_n(A)_1$.
Testing on finitely supported vectors and then using density proves
the same bound for the full commutator. Taking the supremum over
$\varphi$ proves \eqref{eq:commutant-cb}.

For $h\in A'_1$, approximate each contraction in $M_n(B)$ by one in
$M_n(A)$, with arbitrary slack above $\gamma$. It follows that
$\norm{\ad_h|_B}_{\cb}\le2\gamma$.
The first assertion gives distance at most $2\gamma$ to $B'$, and
normalization of approximants gives distance at most $4\gamma$ to
$B'_1$. Reverse the algebras to finish.
\end{proof}

\begin{lemma}[Close representations]\label{lem:paired}
Let $\rho,\theta:P\to\B(H)$ be normal unital $*$-representations of a
von Neumann algebra. If $\norm{\rho-\theta}_{\cb}\le e<1/4$, there
is a unitary $u\in\B(H)$ with
\[
 u\rho(x)u^*=\theta(x)\quad(x\in P),\qquad \norm{u-I}\le4e.
\]
\end{lemma}

\begin{proof}
If $e=0$, take $u=I$. Otherwise let
\[
 \pi(x)=\begin{pmatrix}\rho(x)&0\\0&\theta(x)\end{pmatrix},
 \qquad S=\begin{pmatrix}0&I\\I&0\end{pmatrix}.
\]
The von Neumann image $\pi(P)$ has the quotient matrix norms of $P$.
The commutator of $S$ with $\pi$ has cb norm at most $e$, hence the
same is true on $\pi(P)$ by quotient lifts with arbitrarily small
norm slack. By \Cref{lem:commutant-cb} choose $D\in\pi(P)'$ with
$\norm{D-S}<2e$. Its lower off-diagonal block $T$ satisfies
\[
 T\rho(x)=\theta(x)T,\qquad \norm{T-I}<2e<1.
\]
Thus $T$ is invertible and $T^*T$ commutes with $\rho(P)$.
Its polar unitary $u=T|T|^{-1}$ intertwines the representations.
The singular values of $T$ lie between $1-r$ and $1+r$, where
$r=\norm{T-I}$, so
\[
 \norm{u-I}\le\norm{u-T}+r=\norm{I-|T|}+r\le2r<4e.
\]
\end{proof}

\subsection{The completely positive criterion}

The norm relevant to perturbing a product is the completely bounded
norm on the Haagerup tensor product. We recall only the form needed
here. For a bilinear map $b:P\times P\to P$, it is the least $K$ such
that, at all rectangular sizes,
\begin{equation}\label{eq:haagerup-products}
 \left\|\left[\sum_{j=1}^r b(x_{ij},y_{jk})\right]_{i,k}\right\|
 \le K\norm X\norm Y,
 \quad X\in M_{p,r}(P),\quad Y\in M_{r,q}(P).
\end{equation}
Indeed the matrix Haagerup norm is the infimum of $\norm X\norm Y$
over such row-column factorizations. We write $\norm b_{\cb}$ for
this norm as well.

Write $m_P(x,y)=xy$ for the original multiplication on $P$.
Correction of multiplication by cohomological methods goes back to
Johnson and Raeburn--Taylor
\cite{Johnson1977,RaeburnTaylor1977}. We use the completely bounded
version supplied by Ricard--Roydor's multiplication perturbation theorem
\cite[Proposition~3.2]{RicardRoydor2014}: on a von Neumann algebra $P$,
an associative product $m$ with
$\norm{m-m_P}_{\cb}\le1/11$ admits a cb linear isomorphism $L$ with
\[
 L(m(x,y))=L(x)L(y),\qquad
 \norm{L-\id}_{\cb}\le10\norm{m-m_P}_{\cb}.
\]
If $m(x,y)^*=m(y^*,x^*)$, $L$ can preserve the involution.
There is no normality or separability assumption on the new product.

\begin{proposition}[UCP criterion]\label{prop:ucp}
Let $P,Q$ be von Neumann algebras, let
$\rho:P\to\B(H)$ be a faithful normal unital representation, and let
$Q\subseteq\B(H)$ have the same identity. Suppose that $F:P\to Q$ is
UCP and
\[
 \norm{F-\rho}_{\cb}\le e<10^{-3},\qquad
 \sup_{y\in Q_1}\dist(y,\rho(P)_1)\le g,\qquad g+e<1.
\]
Then there is a normal unital $*$-isomorphism $\theta:P\to Q$ and a
unitary $u\in\B(H)$ such that
\[
 \norm{\theta-\rho}_{\cb}\le32e,\qquad
 u\rho(x)u^*=\theta(x)\ (x\in P),\qquad \norm{u-I}\le128e.
\]
The map $F$ need not be normal.
\end{proposition}

\begin{proof}
Complete isometry of $\rho$ gives
$\norm{F_n(X)}\ge(1-e)\norm X$ at every matrix size. Thus $F$ is
injective, its range $V$ is norm closed, and
$\norm{F^{-1}:V\to P}_{\cb}\le(1-e)^{-1}$.
For $y\in Q_1$, match $y$ with $\rho(x)$ for $x\in P_1$, allowing
arbitrary slack, to obtain $\dist(y,V)\le g+e<1$.
If $V$ were proper, the quotient map $Q\to Q/V$ would have norm one,
contradicting this bound on $Q_1$. Hence $F$ is onto.

Transport multiplication from $Q$ to $P$:
\[
 m(x,y)=F^{-1}(F(x)F(y)).
\]
It is associative and preserves the involution. For rectangular $X,Y$,
insert $\rho(XY)=\rho(X)\rho(Y)$ to see that
\[
 \norm{F(XY)-F(X)F(Y)}\le3e\norm X\norm Y.
\]
Here the first difference is bounded by $e\norm{XY}$ and the product
difference by $2e\norm X\norm Y$, since $F$ and $\rho$ are complete
contractions. Equation~\eqref{eq:haagerup-products} therefore gives
\[
 q:=\norm{m-m_P}_{\cb}\le\frac{3e}{1-e}<\frac1{11}.
\]
The cited theorem gives an involution-preserving $L$ with
$\norm{L-\id}_{\cb}\le\lambda:=30e/(1-e)<1$.
Then $\theta=FL^{-1}$ is a bijective $*$-homomorphism, and
\[
 \norm{\theta-\rho}_{\cb}
 \le e+\norm{L^{-1}-\id}_{\cb}
 \le e+\frac{\lambda}{1-\lambda}
 = e+\frac{30e}{1-31e}\le32e.
\]
Bijective multiplicativity makes $\theta$ unital. A $*$-isomorphism
between von Neumann algebras is an order isomorphism on their
selfadjoint parts, so it preserves bounded increasing suprema and is
normal. Apply \Cref{lem:paired}, since $32e<1/4$.
When $e=0$, the preceding range argument gives $Q=\rho(P)$ directly.
\end{proof}

We will use this criterion both to descend from a tensor product and
to replace a pointwise weak limit by a spatial isomorphism. Its
quantitative conclusion depends only on the cb error of the map;
the ordinary reverse gap supplies surjectivity.

\subsection{Tensor products and averaging}

\begin{lemma}\label{lem:tensor-gap}
Let $A,B\subseteq\B(H)$ be unital von Neumann algebras with
$\dcb(A,B)\le\gamma$. If $K$ is separable, then
\[
 \dcb(A\barotimes\B(K),B\barotimes\B(K))\le\gamma.
\]
If $D\subseteq\B(K)$ is a von Neumann algebra whose commutant is the
strong closure of increasing finite-dimensional unital algebras, then
\[
 \dcb(A\barotimes D,B\barotimes D)\le\gamma.
\]
\end{lemma}

\begin{proof}
Compress a contraction of $M_n(A\barotimes\B(K))$ by increasing
finite-rank projections on $K$. In a corner of rank $r$, choose an
opposite contraction in $M_{nr}(B)$ within $\gamma$ plus arbitrarily
small slack, and extend it by zero on the complement. The compressions
of the original contraction converge strongly to it. A bounded WOT
cluster of the opposite contractions is in $M_n(B\barotimes\B(K))_1$
and is within $\gamma$, by WOT closedness of the norm bound, choosing
the approximation slack to tend to zero along the compressions. Reverse
the argument and take the supremum over $n$.

Now start with a contraction in $M_n(A\barotimes D)$. The first
assertion supplies a nearby contraction in
$M_n(B\barotimes\B(K))$. Average it by conjugation over the unitary
group of each finite-dimensional stage of $D'$. The input commutes
with these groups, so the bound is preserved. A WOT cluster of the
averages commutes with $D'$ and still lies in
$M_n(B\barotimes\B(K))$. The von Neumann tensor commutation theorem
\cite[Theorem~III.4.5.8 and Corollary~III.4.5.9, pp.~319--320]{BlackadarOperatorAlgebras}
places it in $M_n(B\barotimes D)$. Reverse and amplify as before.
\end{proof}

\begin{lemma}[Invariant operator means]\label{lem:invariant-mean}
Let $G$ be a discrete amenable group and $m$ a left invariant mean on
$\ell^\infty(G)$: thus $m$ is a positive unital linear functional with
$m(f(h^{-1}\,\cdot))=m(f)$ for every $h\in G$.
For a bounded family $T_g\in\B(K)$, define
\[
 \ip{(m_gT_g)\xi}{\eta}=m_g\ip{T_g\xi}{\eta}.
\]
This averaging map from bounded operator families to $\B(K)$ is UCP,
commutes with multiplication on either side by a fixed operator, and
preserves weak-operator closed linear subspaces. If
$\norm{T_g-T}\le a$ for every $g$, then
$\norm{m_gT_g-T}\le a$, also at rectangular matrix levels.
No measurability of the family is required.

For a unitary representation $g\mapsto V_g$ on $K$, the map
\[
 E_V(T)=m_gV_gTV_g^*\qquad(T\in\B(K))
\]
is a UCP projection onto $\{V_g:g\in G\}'$. It preserves every von
Neumann algebra invariant under these conjugations.
More generally, for an action by automorphisms $\alpha_g$ on a von
Neumann algebra $P$, the formula
\[
 \varphi(E(x))=m_g\varphi(\alpha_g(x))\qquad(\varphi\in P_*)
\]
defines a UCP map onto $P^G$ fixing that algebra. These averaging maps
need not be normal.
\end{lemma}

\begin{proof}
The coefficient formula defines a bounded sesquilinear form and hence
an operator. Positivity at every matrix size and unitality follow by
applying the positive scalar mean to matrix quadratic forms; the domain
has the pointwise operations on bounded operator families. Fixed left
and right factors pass through the coefficient formula. Testing the
linear equations defining a weak-operator closed subspace proves range
preservation. Testing vector coefficients at each rectangular size
gives the norm bound. For the action average, left invariance makes
the range invariant, while averaging fixes each invariant operator.
An invariant algebra is preserved by the preceding range property.
Predual duality gives the same argument for the general action on $P$,
including preservation of matrix norm bounds.
\end{proof}

For $G=\R$ with the discrete topology, an invariant mean is obtained
as a weak-star cluster point of uniform averages on finite boxes in
finitely generated subgroups. Direct the choices by a finite set of
translations and an error tolerance. Each subgroup is free abelian,
and the relative boundary of its boxes tends to zero for every
prescribed translation; the cluster is therefore positive, unital,
and translation invariant.

The boundary means in \Cref{sec:boundary} require a separate
construction: their equivariance involves a nonsingular action and
their module coefficients vary over the boundary space.
